
% La numeración sigue contando desde la portada, pero no se imprime aquí
\thispagestyle{empty}
\coverheader{%
  PONTIFICIA UNIVERSIDAD CATÓLICA DE CHILE\\
  ESCUELA DE INGENIERÍA
}

% Los bloques se separan con glue elástica (\vfill) en vez de medidas fijas: el
% sobrante de la página se reparte por igual, de modo que el conjunto
% título/autor queda centrado entre el encabezado y el bloque de la carrera.
% Ojo: tiene que ser \vfill (orden `fill') y no \stretch/\vfil, porque \newpage
% agrega su propio \vfil y se llevaría una parte del sobrante.
\vfill

\begin{coverbody}
  % El \par va DENTRO del grupo de \LARGE/\large: si se cierra antes, el
  % párrafo se rompe con el interlineado de 12pt y las líneas se solapan.
  \begin{center}
    {\LARGE\textbf{(TITULO TRABAJO)}\par}

    \vspace{2cm}
    {\large\textbf{(NOMBRE COMPLETO)}\par}
  \end{center}
\end{coverbody}

\vfill

\begin{coverbody}
  Trabajo de Título para optar al título de Ingeniero Civil en ...

  \vspace{1cm}

  Profesor Guía:\\
  \textbf{(NOMBRE PROFE)}
\end{coverbody}

\vfill

% La guía pide solo el año: "Santiago de Chile, (año)"
\begin{coverbody}
  Santiago de Chile, [AÑO]
\end{coverbody}
% Deja la fecha un poco por encima del punto donde termina la línea
% vertical (a 2,5 cm del borde inferior), como en la plantilla oficial.
\vspace{0.5cm}
\newpage
